\documentclass[11pt]{article}
\usepackage[margin=1.1in]{geometry}
\usepackage{amsmath,amssymb,amsthm,booktabs}
\usepackage{float}
\usepackage[colorlinks=true,linkcolor=blue,citecolor=blue,urlcolor=blue]{hyperref}

\newtheorem{theorem}{Theorem}[section]
\newtheorem{proposition}[theorem]{Proposition}
\newtheorem{lemma}[theorem]{Lemma}
\newtheorem{corollary}[theorem]{Corollary}
\newtheorem{conjecture}[theorem]{Conjecture}
\theoremstyle{remark}
\newtheorem{remark}[theorem]{Remark}

\newcommand{\lam}{\lambda}
\DeclareMathOperator{\conj}{conj}

\title{New values of the largest Littlewood--Richardson coefficient\\
via iterated Okounkov compression}
\author{Andre Ruiz Loera\thanks{Affiliation: Department of Mathematics, University of Illinois Urbana-Champaign. Email: \texttt{andrer3@illinois.edu}.}}
\date{\today}

\begin{document}
\maketitle

\begin{abstract}
Let $C(n)=\max c^{\lam}_{\mu\nu}$ over all triples with $|\mu|+|\nu|=|\lam|=n$
(OEIS A387851). Pak and Soskin recently computed $C(n)$ for $24\le n\le 45$ by searching
all partition pairs at every split $k$,
$\sum_{k\le n/2}p(k)p(n-k)$ pairs in all; the last degree took 20 hours on 32
cluster cores. They conjectured that every maximum sits on a nested pair
$\mu\subseteq\nu\subseteq\lam$ with $\nu/\mu$ a disjoint union of squares, uniquely so up
to symmetry. Our starting point is the Okounkov inequality, proved
by Lam, Postnikov and Pylyavskyy, which bounds $c^{\lam}_{\mu\nu}$ by the coefficient of
the pair obtained by halving
the sum vector $w=\mu+\nu$. Together with conjugation symmetry it gives a map $T$ on
partitions of $n$ that never decreases the largest coefficient available at a sum vector.
Pak and Soskin's structure theorem applies the Lam--Postnikov--Pylyavskyy inequality
once, then the averaging step twice with a conjugation in between; we simply keep
iterating until the orbit closes, and that turns the inequality into a search space
reduction. The $p(n)$ sum vectors collapse onto a much smaller set of terminal vectors,
expanding the terminal pairs alone already determines $C(n)$, and a far
smaller scan over the preimages of the winning terminal vectors provably recovers every
maximizing triple. The resulting
algorithm reproduces all 22 reported values $C(24),\dots,C(45)$ on a single laptop
core, taking 34 minutes end to end at $n=45$, and extends the sequence by eight terms,
\[
C(46),\dots,C(53) \;=\; 5932,\ 8481,\ 9926,\ 12828,\ 15762,\ 18570,\ 22982,\ 31264 .
\]
The accompanying maximizer search is exhaustive rather than heuristic, and both parts of
the Pak--Soskin conjecture are verified at each new value: every maximal triple is
square-union nested, and the maximizing triple is unique up to interchange of $\mu$ and
$\nu$ and simultaneous conjugation.
\end{abstract}

\smallskip
\noindent\textbf{2020 Mathematics Subject Classification.}\\
Primary 05E05; Secondary 05A17, 05-08, 20C30.

\smallskip
\noindent\textbf{Keywords.}
Largest Littlewood--Richardson coefficient, Okounkov inequality, Schur positivity,
integer partitions, computational combinatorics, OEIS A387851.

\section{Introduction}

For partitions $\mu,\nu,\lam$ with $|\mu|+|\nu|=|\lam|$, the Littlewood--Richardson
coefficient $c^{\lam}_{\mu\nu}$ is the multiplicity of the Schur function $s_\lam$ in
$s_\mu s_\nu$. Following a question of Stanley \cite[Exc.~82(c)]{Sta23}, set
\[
C(n)\;=\;\max\bigl\{\,c^{\lam}_{\mu\nu}\;:\;|\mu|+|\nu|=|\lam|=n\,\bigr\},
\]
the largest Littlewood--Richardson coefficient at degree $n$; the sequence is
OEIS~A387851 \cite{OEIS}. Stanley showed that the maximum is attained at
$|\mu|\sim n/2$ and that $C(n)=2^{n/2-O(\sqrt{n})}$, and asked for the asymptotic limit
shape; Pak, Panova and Yeliussizov \cite{PPY19} answered that question, proving that all
three partitions have Vershik--Kerov--Logan--Shepp shape, and tabulated $C(n)$ for
$n\le 23$. Pak and Soskin \cite{PS26} computed the full table of
$C(n,k)=\max\{c^{\lam}_{\mu\nu}:|\mu|=k,\ |\lam|=n\}$ for $24\le n\le 45$ by direct
search over triples of partitions, reported that $n=45$ required 20 hours on 32 cores,
and stated that going beyond appeared infeasible by their method. Based on the data they
conjectured:

\begin{conjecture}[{Pak--Soskin \cite[Conj.~1.1]{PS26}}]\label{conj:ps}
For $n\ge 16$, every triple attaining $C(n)$ satisfies
$\mu\subseteq\nu\subseteq\lam$ (up to swapping $\mu,\nu$) with $\nu/\mu$ a disjoint
union of square shapes; for $n\ge 28$ the maximal triple is unique up to
interchange of $\mu$ and $\nu$ and simultaneous conjugation.
\end{conjecture}

Here and throughout,
\emph{conjugation} conjugates all three partitions simultaneously, and a triple is
\emph{square-union nested} if, after possibly interchanging $\mu$ and $\nu$, it
satisfies $\mu\subseteq\nu\subseteq\lam$ and every edge-connected component of
$\nu/\mu$ is a full $c\times c$ square.

This note extends the computation by eight terms and verifies
Conjecture~\ref{conj:ps} at each of them.

\begin{theorem}[computer-assisted]\label{thm:main}
$C(46),\dots,C(53)=5932,\ 8481,\ 9926,\ 12828,\ 15762,\ 18570,\ 22982,\ 31264$. At
each of these $n$, every maximal triple is square-union nested and the maximizing triple
is unique up
to interchange of $\mu$ and $\nu$ and simultaneous conjugation, so both parts of
Conjecture~\ref{conj:ps} hold for $46\le n\le 53$.
\end{theorem}

Section~\ref{sec:method} proves that the search underlying Theorem~\ref{thm:main} is
exhaustive, and Section~\ref{sec:validation} reports the gates
the implementation was put through, including an exhaustive brute force comparison at
every degree $2\le n\le 28$.

The engine is an inequality conjectured by Okounkov \cite{Oko97} and proved in full
generality by Lam, Postnikov and Pylyavskyy \cite[Thm.~11]{LPP07}. The inequality is
theirs; what is done here is to iterate it. In the proof of their Theorem~1.2,
\cite{PS26} first apply the Lam--Postnikov--Pylyavskyy inequality once and then the
Okounkov averaging step twice, with a conjugation in between, to exhibit \emph{some}
maximizer of a prescribed shape. Applying the averaging step repeatedly, until the induced map
on sum vectors closes up on itself, turns the inequality into a compression map that
shrinks the search space by more than three orders of magnitude and, more importantly,
leaves a provably complete list of candidates for the maximizers. That is the entire
content of the method, described next.

The Schur positivity landscape around this inequality has moved quickly in 2026.
Chan, Chen, Pak and Soskin \cite{CCPS26} generalized the
Ahlswede--Daykin inequality to a Schur positive ADS inequality that contains the
Lam--Postnikov--Pylyavskyy inequality as a special case. Speyer \cite{Spe26} proved
the log-concavity conjecture of Lam, Postnikov and Pylyavskyy, Le and Nguyen
\cite{LN26} have since extended that result to skew Schur functions, and,
closest to home,
Chan and Pak \cite{CP26} determined equality conditions for correlation inequalities,
the Okounkov inequality among them. None of that machinery is needed below: the
argument uses the coefficientwise inequality (O) and nothing else. What is new here
is the decision to iterate the averaging and conjugation step on sum vectors, and
to treat the finite dynamics that fall out as an exact computational compression.

\section{The compression map and why the search is complete}\label{sec:method}

Partitions are padded with trailing zeros where needed, and sums such as $\mu+\nu$ are
entrywise. Write $\lfloor w/2\rfloor$ and $\lceil w/2\rceil$ for the entrywise floor and ceiling
halves of an integer vector $w$, and $\conj$ for partition conjugation. Two classical
facts are used, and nothing else:
\begin{enumerate}
\item[(O)] (Okounkov inequality \cite{Oko97,LPP07}) For all partitions $\mu,\nu$ and
every $\lam$,
\[
c^{\lam}_{\mu\nu}\;\le\;c^{\lam}_{\;\lfloor(\mu+\nu)/2\rfloor,\;\lceil(\mu+\nu)/2\rceil}.
\]
\item[(C)] (Conjugation symmetry) $c^{\lam}_{\mu\nu}=c^{\conj\lam}_{\conj\mu,\conj\nu}$.
\end{enumerate}
The right-hand side of (O) depends on $(\mu,\nu)$ only through the sum vector
$w=\mu+\nu$, which is itself a partition of $n$. For a partition $w$ of $n$ put
\[
M(w)\;=\;\max_{\lam\,\vdash\,n}\ c^{\lam}_{\;\lfloor w/2\rfloor,\;\lceil w/2\rceil},
\qquad
T(w)\;=\;\conj\bigl(\lfloor w/2\rfloor\bigr)+\conj\bigl(\lceil w/2\rceil\bigr) .
\]
Thus $T$ halves $w$, conjugates both halves, and records the new sum vector; $T(w)$ is
again a partition of $n$.

\begin{lemma}\label{lem:mono}
Let $w$ be a partition of $n$. Then
\begin{enumerate}
\item[(a)] $M(w)=\max\bigl\{c^{\lam}_{\mu\nu}\ :\ \mu+\nu=w\bigr\}$, the maximum being
over all pairs of partitions with sum vector $w$ and all $\lam$;
\item[(b)] $M(w)\le M\bigl(T(w)\bigr)$.
\end{enumerate}
\end{lemma}

\begin{proof}
(a) The pair $\bigl(\lfloor w/2\rfloor,\lceil w/2\rceil\bigr)$ has sum vector $w$, so
$M(w)$ is at most the right-hand side. Conversely, if $\mu+\nu=w$ then (O) gives
$c^{\lam}_{\mu\nu}\le c^{\lam}_{\lfloor w/2\rfloor,\lceil w/2\rceil}\le M(w)$ for every
$\lam$.

(b) Choose $\lam$ with $c^{\lam}_{\lfloor w/2\rfloor,\lceil w/2\rceil}=M(w)$. By (C),
\[
M(w)\;=\;c^{\conj\lam}_{\;\conj\lfloor w/2\rfloor,\;\conj\lceil w/2\rceil}.
\]
The pair $\bigl(\conj\lfloor w/2\rfloor,\conj\lceil w/2\rceil\bigr)$ is a pair of
partitions with sum vector $T(w)$, by the definition of $T$. Applying part~(a) at
$T(w)$ bounds the right-hand side by $M(T(w))$.
\end{proof}

Since there are only $p(n)$ partitions of $n$, the sequence $w,T(w),T^2(w),\dots$ must
repeat. Let $k(w)\ge 1$ be least with $T^{k(w)}(w)\in\{w,T(w),\dots,T^{k(w)-1}(w)\}$ and
put
\[
w^\ast\;=\;T^{k(w)}(w),
\]
the first repeated iterate. By construction $w^\ast$ recurs under $T$, so it is a
periodic point of $T$; we call it the \emph{terminal vector} of $w$, and call
$\bigl(\lfloor w^\ast/2\rfloor,\lceil w^\ast/2\rceil\bigr)$ the \emph{terminal pair}.
Iterating Lemma~\ref{lem:mono}(b) along $w,T(w),\dots,T^{k(w)}(w)$ gives
\begin{equation}\label{eq:chain}
M(w)\;\le\;M(w^\ast)\qquad\text{for every }w\vdash n .
\end{equation}

\begin{proposition}\label{prop:complete}
Let $n\ge 1$. Then
\begin{enumerate}
\item[(a)] $\displaystyle C(n)=\max_{w\,\vdash\,n} M(w^\ast)$, the maximum being over the
terminal vectors alone;
\item[(b)] if $c^{\lam}_{\mu\nu}=C(n)$ and $w=\mu+\nu$, then $M(w^\ast)=C(n)$.
\end{enumerate}
\end{proposition}

\begin{proof}
By Lemma~\ref{lem:mono}(a), $C(n)=\max_{w\vdash n}M(w)$.

(a) Each $w^\ast$ is itself a partition of $n$, so $\max_w M(w^\ast)\le\max_w
M(w)=C(n)$; and \eqref{eq:chain} gives $M(w)\le M(w^\ast)\le \max_{w}M(w^\ast)$ for
every $w$, hence $C(n)\le\max_w M(w^\ast)$.

(b) Lemma~\ref{lem:mono}(a) gives $M(w)\ge c^{\lam}_{\mu\nu}=C(n)$, and \eqref{eq:chain}
gives $M(w^\ast)\ge M(w)\ge C(n)$, while (a) gives $M(w^\ast)\le C(n)$.
\end{proof}

Call a terminal vector \emph{winning} if $M(w^\ast)=C(n)$, and call a partition $w$ of
$n$ a \emph{suspect} if $w^\ast$ is winning.

\begin{corollary}\label{cor:exhaustive}
Every triple $(\lam,\mu,\nu)$ with $c^{\lam}_{\mu\nu}=C(n)$ has $\mu+\nu$ a suspect.
Consequently, expanding $s_\mu s_\nu$ for every decomposition $\mu+\nu=w$ of every
suspect $w$ recovers \emph{every} maximal triple, and the resulting verification of
Conjecture~\ref{conj:ps} is exhaustive, not heuristic.
\end{corollary}

\begin{proof}
Immediate from Proposition~\ref{prop:complete}(b) and the definitions.
\end{proof}

\begin{remark}[the iteration does not converge to a fixed point]\label{rem:cycles}
Nothing above asserts that $w^\ast$ is fixed by $T$, and in general it is not. At six of
the eight degrees treated below the sum vector of the maximizing pair lies on a genuine
$2$-cycle of $T$, and only at $n=46$ and $n=53$ is it a fixed point. Only
Lemma~\ref{lem:mono}(b) and the finiteness of the state space are used, so the argument
is indifferent to which. Two directly verified facts explain two of the numbers reported
below. First, at every degree $46\le n\le 53$ the set of terminal vectors was verified
to be exactly the set of partitions of $n$ whose odd parts are pairwise distinct, so the
terminal class counts in Table~\ref{tab:results} are values of OEIS~A006950 \cite{OEIS}.
Second, see Lemma~\ref{lem:cycleconst}. The dynamics of $T$ in their own right are taken
up in a separate note; nothing in the present one depends on them.
\end{remark}

\begin{lemma}\label{lem:cycleconst}
$M$ is constant on every cycle of $T$. Hence the set of winning terminal vectors is a
union of $T$-cycles, and its cardinality is a sum of cycle lengths.
\end{lemma}

\begin{proof}
If $T^m(v)=v$ then Lemma~\ref{lem:mono}(b) gives
$M(v)\le M(T(v))\le\dots\le M(T^m(v))=M(v)$, so all the inequalities are equalities.
\end{proof}

At each degree $46\le n\le 53$ the winning set turns out to be a \emph{single} cycle:
a $2$-cycle except at $n=46$ and $n=53$, where it is a fixed point. This is the column
``classes attaining'' of Table~\ref{tab:gap}, whose entries $1$ and $2$ are therefore an
artefact of the cycle structure and not a failure of uniqueness. At the six degrees with
a $2$-cycle, the two winning terminal pairs are conjugates of one another, so they carry
the same maximal triples up to conjugation.

\subsection*{The computation}

\begin{enumerate}
\item \textbf{Compress.} For each $w\vdash n$ compute $w^\ast$ by iterating $T$
(memoized); record the fibers $w\mapsto w^\ast$. Cost: seconds.
\item \textbf{Expand terminal pairs.} For each distinct terminal vector $w^\ast$, expand
$s_{\lfloor w^\ast/2\rfloor}\, s_{\lceil w^\ast/2\rceil}$ with \texttt{lrcalc}
\cite{lrcalc} and take the largest coefficient; this is $M(w^\ast)$. By
Proposition~\ref{prop:complete}(a) the maximum over terminal vectors is $C(n)$. This pass
dominates the runtime.
\item \textbf{Exhaust the suspects.} Expand $s_\mu s_\nu$ for every decomposition
$\mu+\nu=w$ of every suspect $w$. By Corollary~\ref{cor:exhaustive} this recovers every
maximal triple $(\lam,\mu,\nu)$.
\end{enumerate}
The $p(n)$ sum vectors collapse onto a much smaller set of terminal vectors: $12\,306$ at
$n=46$ growing to $32\,490$ at $n=53$, against $2.1\cdot 10^{7}$ pairs $(\mu,\nu)$ with
$|\mu|\le|\nu|$ at $n=46$ growing to $1.0\cdot 10^{8}$ at $n=53$. In pass~3 the suspects
are few: tens of sum vectors and $10^3$ to $10^4$ pairs at every degree computed.

Each maximal triple is then tested for the square-union property (each edge-connected
component of the skew shape $\nu/\mu$ is a full $c\times c$ square, the reading of
\cite{PS26} matching their $n=40$ example, in which four unit boxes touching only
diagonally count as four squares), and the list is collapsed under the interchange and
conjugation action of Conjecture~\ref{conj:ps} to count orbits.

\section{Validation}\label{sec:validation}

The method was gated five ways before any new value was believed.

\medskip
\noindent\textbf{(V1) Exhaustive brute force at every degree $2\le n\le 28$.}
For each such $n$ we computed $C(n)$ \emph{and the complete list of maximal triples} by
direct search over every pair $(\mu\vdash k,\ \nu\vdash n-k)$ with $k\le n/2$, sharing no
step with the compression, and compared both against the three-pass method. At all $27$
degrees the values agree and the two lists of maximal triples agree \emph{as sets};
so does the verdict on each part of Conjecture~\ref{conj:ps}. This is a direct test of
Corollary~\ref{cor:exhaustive}: had the suspect scan been able to miss a maximizer, the
lists would differ. Script \texttt{bruteforce\_gate.py}, log
\texttt{bruteforce\_gate.log}.

\medskip
\noindent\textbf{(V2) The full reported range.} The implementation reproduces all $22$
values $C(24),\dots,C(45)$ of the appendix of \cite{PS26}, in one run. Those values were
not transcribed: they were parsed out of the \LaTeX{} source of the \texttt{arXiv}
e-print of \cite{PS26}, and the maxima of the parsed $C(n,k)$ grid were checked against
the parsed $C(n)$ row (all 22 rows, no discrepancy). Runtimes at the degrees
\cite{PS26} timed:

\begin{center}
\begin{tabular}{rrrrr}
\toprule
$n$ & $C(n)$ \cite{PS26} & $C(n)$ here & runtime \cite{PS26} & runtime here (1 core) \\
\midrule
25 & 45   & 45   & 28 s (4 cores) & 0.8 s \\
30 & 176  & 176  & 198 s (4 cores) & 4.6 s \\
35 & 423  & 423  & 3\,258 s (4 cores) & 40 s \\
40 & 1\,484 & 1\,484 & 44\,919 s (4 cores) & 234 s \\
41 & 1\,768 & 1\,768 & 5\,609 s (32 cores) & 387 s \\
45 & 4\,323 & 4\,323 & 72\,172 s (32 cores) & 2\,042 s \\
\bottomrule
\end{tabular}
\end{center}

\noindent
The $n=41$ row is the one where \cite{PS26} moved from a laptop to the cluster, which is
why their time drops there. All six of our times come from the single run recorded in
\texttt{revalidate\_24\_45.log}, and all six are totals over the three passes. At
$n=45$ the time splits as $15$ s of compression, $1\,429$ s of terminal expansion and
$598$ s of suspect scanning. Our times are single core, and the machine was carrying
heavy unrelated load throughout, so they are upper bounds on the cost of the method
rather than benchmarks. The same caveat applies to Table~\ref{tab:results}, whose
degrees were run at different times and under different load, which is why $n=51$
there exceeds $n=52$.

\medskip
\noindent\textbf{(V3) External answer keys.} Three exhibits printed in \cite{PS26} were
reproduced independently: their $n=40$ maximal triple $\lam=(9,7,6,5,4,3,2,2,1,1)$,
$\mu=(6,4,3,2,2,1)$, $\nu=(6,5,4,3,2,1,1)$ has $c^{\lam}_{\mu\nu}=1484=C(40)$ with
$\nu/\mu$ four unit squares; their counterexample at $n=11$ to the nestedness conjecture
of \cite{PPY19} has $c^{(5,3,2,1)}_{(4,1),(3,2,1)}=3=C(11)$ with $\mu\not\subseteq\nu$,
and our brute force at $n=11$ independently finds part~1 of Conjecture~\ref{conj:ps}
failing there; and their four maximizers at $n=20$ all evaluate to $18=C(20)$, which our
brute force at $n=20$ independently returns together with exactly four maximal triples.
Finally, in the range $2\le n\le 28$ covered by (V1) the uniqueness assertion of
Conjecture~\ref{conj:ps} fails for the last time at $n=27$ and holds at $n=28$, which is
where \cite{PS26} place the threshold; that threshold was not an input to our
computation.

\medskip
\noindent\textbf{(V4) A second Littlewood--Richardson implementation.} The coefficient of
the maximal triple at each of $n=46,\dots,53$ was recomputed by direct enumeration of
Littlewood--Richardson skew tableaux, in a from-scratch counter sharing no code and no
algorithm with \texttt{lrcalc} (script \texttt{crosscheck.py}). All eight agree, as do
four textbook sanity values. In addition, \texttt{lrcalc} was used to recompute the
coefficient of every one of the fourteen maximal triples found at $46\le n\le 53$, not
only the eight representatives printed below.

\medskip
\noindent\textbf{(V5) A second enumeration of the terminal vectors.} Pass~2 was rerun
from an independent script (\texttt{indep53.py}) which enumerates the terminal vectors
directly, as the partitions of $n$ with pairwise distinct odd parts, instead of by
iterating $T$ over all $p(n)$ sum vectors, and which shares no code with
\texttt{maxlr.py}. It has now been run to completion at \emph{every} degree
$46\le n\le 53$, and it reproduces Table~\ref{tab:results} and Table~\ref{tab:gap}
entry for entry: the value $C(n)$, the number of terminal classes, the number of
classes attaining the maximum, and the largest value strictly below it all agree at
all eight degrees. At $n=46$, for instance, it returns $C(46)=5932$ from the same
$12\,306$ terminal vectors with one class attaining and $5\,862$ below, and at $n=53$
it returns $C(53)=31264$ from the same $32\,490$ terminal vectors with one class
attaining and $27\,046$ below. The two enumerations were also verified to produce
literally the same set of terminal vectors at every $46\le n\le 53$.

\medskip
What none of these gates provides is a second, independent confirmation of the
\emph{upper} bound at the new degrees from a different tool chain: the statement that no
pair at degree $n$ beats the reported value rests, for $46\le n\le 53$, on
Proposition~\ref{prop:complete} together with runs of pass~2 over the terminal vectors,
by two scripts but with the same underlying Littlewood--Richardson library. The lower
bounds, that the reported triples achieve the reported coefficients, are confirmed by two
independent implementations.

\section{Results}\label{sec:results}

\begin{table}[H]
\centering
\begin{tabular}{rrrrrrr}
\toprule
$n$ & $C(n)$ & terminal vectors & suspect pairs & maximal triples & orbits & runtime \\
\midrule
46 & 5\,932  & 12\,306 & 8\,477  & 1 & 1 & 354 s \\
47 & 8\,481  & 14\,193 & 12\,032 & 2 & 1 & 590 s \\
48 & 9\,926  & 16\,357 & 14\,304 & 2 & 1 & 849 s \\
49 & 12\,828 & 18\,803 & 14\,884 & 2 & 1 & 1\,345 s \\
50 & 15\,762 & 21\,581 & 23\,084 & 2 & 1 & 2\,239 s \\
51 & 18\,570 & 24\,766 & 8\,910  & 2 & 1 & 11\,209 s \\
52 & 22\,982 & 28\,394 & 4\,901  & 2 & 1 & 9\,922 s \\
53 & 31\,264 & 32\,490 & 2\,835  & 1 & 1 & 13\,429 s \\
\bottomrule
\end{tabular}
\caption{The eight new values. ``Maximal triples'' counts every triple attaining $C(n)$
with $|\mu|\le|\nu|$; ``orbits'' counts them up to the interchange and conjugation
action of Conjecture~\ref{conj:ps}. Runtimes are a single core of a consumer laptop
(Apple M4, Section~\ref{sec:repro}) which was carrying unrelated load throughout, which
is why $n=51$ exceeds $n=52$.}
\label{tab:results}
\end{table}

Representatives of the unique maximal orbit at each degree:
\begin{align*}
n=46:&\quad \lam=(10,8,7,6,4,4,3,2,1,1),\ \mu=\nu=(7,5,4,3,2,1,1);\\
n=47:&\quad \lam=(10,8,7,6,5,4,3,2,1,1),\ \mu=(7,5,4,3,2,1,1),\ \nu=(7,6,4,3,2,1,1);\\
n=48:&\quad \lam=(10,9,7,6,5,4,3,2,1,1),\ \mu=\nu=(7,6,4,3,2,1,1);\\
n=49:&\quad \lam=(11,9,7,6,5,4,3,2,1,1),\ \mu=(7,6,4,3,2,1,1),\ \nu=(8,6,4,3,2,1,1);\\
n=50:&\quad \lam=(12,9,7,6,5,4,3,2,1,1),\ \mu=\nu=(8,6,4,3,2,1,1);\\
n=51:&\quad \lam=(11,9,7,6,5,4,3,2,2,1,1),\ \mu=(7,5,4,3,2,1,1),\ \nu=(8,6,5,3,2,2,1,1);\\
n=52:&\quad \lam=(11,9,8,6,5,4,3,2,2,1,1),\ \mu=(7,5,4,3,2,1,1),\ \nu=(8,6,5,4,3,2,1);\\
n=53:&\quad \lam=(11,9,8,6,5,4,3,3,2,1,1),\ \mu=(7,5,4,3,2,1,1),\ \nu=(8,6,5,4,3,2,1,1).
\end{align*}
In every case $\nu/\mu$ is a disjoint union of unit squares no two of which share an edge
(at $n=46,48,50$ it is empty, $\mu=\nu$, vacuously such a union), confirming part one of
Conjecture~\ref{conj:ps}; the orbit count $1$ in Table~\ref{tab:results} is part two. The
second maximal triple at the six degrees where there are two is the conjugate of the
first, and is printed in the run logs. The largest consecutive ratio in the extended
table is $C(47)/C(46)\approx 1.430$ and the smallest is $C(48)/C(47)\approx 1.170$, a
spread comfortably inside the range $1.037$ to $1.576$ of the consecutive ratios in
$24\le n\le 45$: the growth stays as irregular as it was.

\subsection*{The gap below the maximum}

\begin{table}[H]
\centering
\begin{tabular}{rrrrr}
\toprule
$n$ & $C(n)$ & terminal vectors attaining & largest value below & ratio \\
\midrule
46 & 5\,932  & 1 & 5\,862  & 1.0119 \\
47 & 8\,481  & 2 & 7\,154  & 1.1855 \\
48 & 9\,926  & 2 & 9\,413  & 1.0545 \\
49 & 12\,828 & 2 & 12\,449 & 1.0304 \\
50 & 15\,762 & 2 & 14\,631 & 1.0773 \\
51 & 18\,570 & 2 & 18\,356 & 1.0117 \\
52 & 22\,982 & 2 & 22\,351 & 1.0282 \\
53 & 31\,264 & 1 & 27\,046 & 1.1560 \\
\bottomrule
\end{tabular}
\caption{The largest value of $M$ over terminal vectors strictly below $C(n)$. By
Lemma~\ref{lem:cycleconst} the ``attaining'' column is a sum of $T$-cycle lengths; here
it is one cycle throughout.}
\label{tab:gap}
\end{table}

The compression framework makes one statistic cheap to record: the largest value
$M(w^\ast)$ over terminal vectors \emph{strictly below} $C(n)$. The ratio column of
Table~\ref{tab:gap} is $C(n)$ divided by that runner-up. Every row comes from
\texttt{top2\_gap.py}, and the $n=53$ row is additionally confirmed by the
independent terminal enumeration of gate (V5), whose per-vector maxima give
the runner-up directly; the two agree on all four quantities
(logs \texttt{top2\_gap\_53.log} and \texttt{indep53.log}). The gap is uniformly positive but far from uniform in size: at
$n=46$ and $n=51$ the runner-up comes within $1.2\%$ of the maximum, while at $n=47$
the winner exceeds the runner-up by $18.6\%$.
Uniqueness up to symmetry holds at every computed $n$ nonetheless, including the two
tight degrees.

We stress what this statistic is not. Proposition~4.1 of \cite{PS26} shows that if $C(n)$
is attained at a \emph{single} $k$ among the values $\{C(n,k):0\le k\le n/2\}$, then
part~one of Conjecture~\ref{conj:ps} holds at $n$ (part one, not the uniqueness part),
and \cite{PS26} then speculate that a runaway growth of $C(n)$ might make that hypothesis
provable asymptotically. The relevant gap there is between the largest and the second
largest entry of the row $\{C(n,k)\}_k$, at most $n/2$ numbers indexed by $k$. That is a
different quantity from the one in Table~\ref{tab:gap}, and the present method does not
deliver it: the compression map does not preserve $\min(|\mu|,|\nu|)$, so the terminal
vectors do not respect $k$ and no row of $C(n,k)$ can be read off them. (The step
$w\mapsto w^\ast$ changes $\min(|\mu|,|\nu|)$ for $74\%$ of the sum vectors at $n=20$
and $82\%$ at $n=30$.) No row of $C(n,k)$ beyond $n=45$ is computed here, and we
are not aware of one anywhere.

\begin{remark}
The method is a search space reduction, not a new upper bound, and the cost stays
superpolynomial. Already the number of partitions of $n$ with distinct odd parts, which
at every computed degree is exactly the number of terminal vectors, grows like
$e^{\Theta(\sqrt n)}$ just as
$p(n)$ does, though on the computed range the terminal count sits between a tenth and an
eighth of $p(n)$. Pass~2 parallelizes
trivially over terminal vectors, and a multi-process or PyPy implementation would push
the table a few terms further; we have not established how far.
\end{remark}

\section{Reproducibility}\label{sec:repro}

All computations were run with CPython 3.9.6 on macOS 15.7.1, on one core of an Apple
M4 laptop (4 performance and 6 efficiency cores, 16 GB RAM), in a single process and a
single thread, with the \texttt{lrcalc} library of
Buch \cite{lrcalc}, version 2.1, through its Python bindings (\texttt{pip install
lrcalc}). Wall-clock times were taken under unrelated load, so read them as upper
bounds; the
machine-independent measure of the work is the terminal-vector and suspect-pair counts
in Table~\ref{tab:results}. The scripts and the raw logs of every run reported here are supplied as
ancillary files with this submission, in the directory \texttt{anc/}:

\begin{center}
\small
\begin{tabular}{ll}
\toprule
file & what it does \\
\midrule
\texttt{maxlr.py} & the three-pass algorithm; \texttt{python3 maxlr.py 46 ... 53} \\
\texttt{crosscheck.py} & independent LR counter, no \texttt{lrcalc}; \texttt{python3 crosscheck.py} \\
\texttt{bruteforce\_gate.py} & gate (V1); \texttt{python3 bruteforce\_gate.py 2 28} \\
\texttt{indep53.py} & independent pass~2; \texttt{python3 indep53.py 53 3} \\
\texttt{top2\_gap.py} & Table~\ref{tab:gap}; \texttt{python3 top2\_gap.py 46 ... 53} \\
\texttt{results\_46\_50.log} & runs behind Table~\ref{tab:results}, $n=46$ to $50$ \\
\texttt{results\_51\_52.log} & runs behind Table~\ref{tab:results}, $n=51,52$ \\
\texttt{results\_53.log} & run behind Table~\ref{tab:results}, $n=53$ \\
\texttt{revalidate\_24\_45.log} & gate (V2), the reported range \\
\texttt{bruteforce\_gate.log} & gate (V1) \\
\texttt{crosscheck.log} & gate (V4) \\
\texttt{indep\_46\_52.log}, \texttt{indep53.log} & gate (V5) \\
\texttt{top2\_gap.log}, \texttt{top2\_gap\_53.log} & Table~\ref{tab:gap} \\
\texttt{README.txt} & the above, with the environment and expected runtimes \\
\bottomrule
\end{tabular}
\end{center}

\noindent
Every value in the tables above appears in one of those logs, and the remaining figures
quoted in the text are arithmetic on them. The range-based scripts take a degree range
on the command line as shown above (\texttt{crosscheck.py} takes no arguments) and
print one line per degree; none of them needs an argument file, a data
file or a network connection.

\section*{Acknowledgements}

Claude (Anthropic) was used for assistance with the implementation, the computations,
and the preparation of this manuscript; every new computational result reported in this paper is an output of the
accompanying code and its committed logs. The author reviewed the proofs, the code and
the logs, and is responsible for the results and conclusions.

\begin{thebibliography}{99}

\bibitem{CCPS26}
S.~H. Chan, H.~Chen, I.~Pak, D.~Soskin,
\emph{Correlation inequalities for Schur positivity},
arXiv:2606.06688 (2026).

\bibitem{CP26}
S.~H. Chan, I.~Pak,
\emph{Equality conditions for correlation inequalities},
arXiv:2607.06275 (2026).

\bibitem{LPP07}
T.~Lam, A.~Postnikov, P.~Pylyavskyy,
\emph{Schur positivity and Schur log-concavity},
Amer. J. Math. \textbf{129} (2007), 1611--1622.

\bibitem{LN26}
T.~Le, S.~Nguyen,
\emph{Skew hives, skew skeps, skew Schur log-concavity},
arXiv:2608.13544 (2026).

\bibitem{lrcalc}
A.~S. Buch,
\emph{The Littlewood--Richardson calculator} (\texttt{lrcalc}), version 2.1,
software, \url{https://sites.math.rutgers.edu/~asbuch/lrcalc/}.

\bibitem{OEIS}
OEIS Foundation Inc.,
\emph{The On-Line Encyclopedia of Integer Sequences}, entries A387851 and A006950,
\url{https://oeis.org}.

\bibitem{Oko97}
A.~Okounkov,
\emph{Log-concavity of multiplicities with applications to characters of $U(\infty)$},
Adv. Math. \textbf{127} (1997), 258--282.

\bibitem{PPY19}
I.~Pak, G.~Panova, D.~Yeliussizov,
\emph{On the largest Kronecker and Littlewood--Richardson coefficients},
J. Combin. Theory Ser. A \textbf{165} (2019), 44--77.

\bibitem{PS26}
I.~Pak, D.~Soskin,
\emph{On the largest Littlewood--Richardson coefficient},
arXiv:2608.03281 (2026).

\bibitem{Spe26}
D.~E. Speyer,
\emph{L-log-concavity and a proof of the conjecture of Lam, Postnikov and Pylyavskyy},
arXiv:2601.05007 (2026).

\bibitem{Sta23}
R.~P. Stanley,
\emph{Supplementary exercises for Enumerative Combinatorics}, vol.~2, Chapter~7,
Exercise~82(c), version of June~4, 2023;
\url{http://www-math.mit.edu/~rstan/ec}.

\end{thebibliography}

\end{document}
